\section{Exploratory Historical Trajectories}
\label{sec:pilot}
We separately analyze three prior runs configured around approximately
100 agents: an adventuring guild (25 rounds), an underground black-market
auction (20), and a storm-struck luxury cruise (20). They predate the current
hypotheses and differ in configuration and social context. We reconstruct
undirected actor--target networks, route entropy, and sums of directed
relationship changes using the same metric definitions.

\begin{table*}[t]
\centering\small
\setlength{\tabcolsep}{5pt}
\begin{tabularx}{\textwidth}{@{}lrrrrrrY@{}}
\toprule
World & Rounds & Events & IDs & Dyads & Repeat & Largest comp. & Dominant state changes \\
\midrule
Guild & 25 & 444 & 103 & 282 & 36.5\% & 11 & Trust $+3473$; affection $+1745$ \\
Black market & 20 & 299 & 99 & 229 & 23.4\% & 13 & Influence/fear $+1194$ \\
Storm cruise & 20 & 309 & 100 & 235 & 23.9\% & 13 & Trust $+2576$; affection $+961$ \\
\bottomrule
\end{tabularx}
\caption{Exploratory statistics from 1,052 historical events. IDs are distinct
logged actor/target identifiers, not a fixed population denominator; the guild
record contains 103. Relationship totals are simulated state changes.}
\label{tab:pilot}
\end{table*}

Between 23.4\% and 36.5\% of events revisit an existing dyad, while each world
contains 229--282 unique dyads (Table~\ref{tab:pilot}). Route entropies are
2.889, 2.619, and 2.524 bits for guild, black market, and cruise, respectively.
Because the repertoires differ, these values describe each trace's variety;
they are not directly comparable measures of social intelligence.

The largest components contain only 11--13 participants despite many
observed IDs. Figure~\ref{fig:pilot-dynamics} shows dyads continuing to
accumulate after component sizes stabilize. This is consistent with additional
ties forming predominantly within local clusters. Spatial encounter and
activation policies can generate such structure, so its social interpretation
must remain conditional on those mechanisms.

\begin{figure*}[t]
\centering
\includegraphics[width=0.96\textwidth]{figures/pilot_social_dynamics.pdf}
\caption{\paneltitle{Longer traces distinguish activity from connectivity.}
Unique dyads continue to accumulate, while the largest observed components
plateau early. These historical runs motivate testing transmission across
local clusters; they do not identify its causal mechanism.}
\label{fig:pilot-dynamics}
\end{figure*}

In the black market, influence/fear ($+1194$) exceeds trust ($+179$) and
affection ($+9$); the guild and cruise instead accumulate mostly trust and
affection. The pattern is compatible with each scenario's institutions, but
prompts, action definitions, schedules, and run histories are confounded.
It motivates a controlled question rather than answering one. In particular,
warning diffusion and collective rescue need separate outcomes: receiving
information need not produce action \cite{centola}. Appendix~\ref{app:intervention}
specifies a paired memory intervention to examine that distinction.
